\documentclass{article}
\usepackage[a4paper, total={6in, 8in}]{geometry}

\usepackage{amsmath}
\usepackage{bm}
\usepackage{amsfonts}
\usepackage{blkarray}
\usepackage{graphicx}
\usepackage{subcaption}
\usepackage{placeins}
\usepackage{mathtools}
\usepackage{multirow}
\usepackage{algorithm}
\usepackage{algpseudocode}
\usepackage{tikz}
\usepackage{nicematrix}
\usepackage{xcolor}

\title{Structure-exploiting LU decomposition for implicit integrators in the Riccati recursion}
\author{Louis Callens}

\begin{document}
\maketitle
\section{Introduction}
In the Riccati recursion, states (and dynamics lagrangian multipliers) are eliminated using the dynamics equation
\begin{equation}
    f(x_k, u_k) - x_{k+1} = 0
\end{equation}
of which the linearization that appears in the constraint jacobian is
\begin{equation}
    B_k u_k + A_k x_k - x_{k+1} = 0
\end{equation}

Controls (and stage-wise equality lagrnagian multipliers) are eliminated partially using the equality constraints by decomposing the stage-wise equality constraint jacobian, given by
\begin{equation}
    H_u = P_l T_l \begin{bmatrix}\bm{I}_\rho & \\ & \bm{0} \end{bmatrix} T_r P_r
\end{equation}

In the case of implicit integrators, the dynamics equation is given by
\begin{equation}
    f(x_k, u_k, x_{k+1}) = 0
\end{equation}
which can be implemented by introducing the controls $z_k \in \mathbb{R}^{n_x}$ and enforcing
\begin{align}
    z_k - x_{k+1} &= 0\\
    f(x_k, u_k, z_k) &= 0
\end{align}
which leads to a specific structure in the stage-wise equality constraint jacobian, since the original equality do not depend on $z_k$
\begin{align}
    H_u &=
    \begin{bmatrix}
        \nabla_{u_k} g_\mathrm{eq}(x_k, u_k) & \nabla_{z_k} g_\mathrm{eq}(x_k, u_k)\\
        \nabla_{u_k} f(x_k, u_k, z_k) & \nabla_{z_k} f(x_k, u_k, z_k)
    \end{bmatrix}
    =
    \begin{bmatrix} M_1 & \\M_2 & M_3 \end{bmatrix}
    \label{eq:Hu}
\end{align}
where $M_1 \in \mathbb{R}^{n_g \times n_u}$, $M_2 \in \mathbb{R}^{n_x \times n_u}$, and $M_3 \in \mathbb{R}^{n_x \times n_x}$.

Note that instead of a lower triangular structre, an upper triangulare structure for $H_u$ is also possible, in which case we write
\begin{align}
    H_u &=
    \begin{bmatrix}
        \nabla_{z_k} f(x_k, u_k, z_k) & \nabla_{u_k} f(x_k, u_k, z_k)\\
        \nabla_{z_k} g_\mathrm{eq}(x_k, u_k) & \nabla_{u_k} g_\mathrm{eq}(x_k, u_k)
    \end{bmatrix}
    =
    \begin{bmatrix} M_1 & M_2\\ & M_3 \end{bmatrix}
    \label{eq:Hu-upper}
\end{align}
where $M_1 \in \mathbb{R}^{n_x \times n_x}$, $M_2 \in \mathbb{R}^{n_x \times n_u}$, and $M_3 \in \mathbb{R}^{n_g \times n_u}$.

Figure \ref{fig:lu_time_spent} shows the percentage of computation time of the Ricatti recursion spent in computing the LU factorization of the matrix shown in \eqref{eq:Hu} as a function of $n_x$, $n_u$, and $n_g$. Especially if $n_g$ and $n_u$ are large, a significant percentage of time is spent in the LU decomposition, motivating the need to perform this decomposition more efficiently. For small values of $n_x$, the decomposition is less significant.

\begin{figure}
    \centering
    \begin{subfigure}{0.3\linewidth}
        \includegraphics[width=\linewidth]{../figures/scaling_2d_lu_time_spent_ng_nx.png}        
    \end{subfigure}
    \begin{subfigure}{0.3\linewidth}
        \includegraphics[width=\linewidth]{../figures/scaling_2d_lu_time_spent_ng_nu.png}        
    \end{subfigure}
    \begin{subfigure}{0.3\linewidth}
        \includegraphics[width=\linewidth]{../figures/scaling_2d_lu_time_spent_nx_nu.png}        
    \end{subfigure}
    \caption{Percentage of computation time of the Ricatti recursion spent in computing LU factorization of the matrix shown in \eqref{eq:Hu} as a function of $n_x$, $n_u$, and $n_g$. Gray colored areas indicate no data available in the performed benchmark.}
    \label{fig:lu_time_spent}
\end{figure}





\section{Blocked Lower Triangular LU decomposition}
The structure in the matrix given by \eqref{eq:Hu} can be exploited by performing the decomposition more efficiently. More precisely, we wish to compute
\begin{equation}
    P H_u Q = L U
\end{equation}
Suppose $M_1$ is decomposed as 
\begin{align}
    P_1 M_1 Q_1 &= L_1 U_1
\end{align}
where
\begin{align}
    L_1 &= \begin{bmatrix} K_1 & \\ K_2 & \bm{I}_{n_g-\rho} \end{bmatrix} & 
    &\text{and} &
    U_1 &= \begin{bmatrix} V_1 & V_2 \\ & \bm{0}_{n_g-\rho \times n_u-\rho} \end{bmatrix}
\end{align}
Define 
\begin{align}
    \tilde{L} &\coloneqq \begin{bmatrix} K_1 & & \\ K_2 & \bm{I}_{n_g-\rho} & \\ & & \bm{I}_{n_x} \end{bmatrix}
    & &\text{such that} &
    \tilde{L}^{-1} &= \begin{bmatrix} K_1^{-1} & & \\ -K_2 K_1^{-1} & \bm{I}_{n_g-\rho} & \\ & & \bm{I}_{n_x} \end{bmatrix}
\end{align}
which leads to
\begin{equation}
    \tilde{L}^{-1} 
    \begin{bmatrix} P_1 & \\ & \bm{I}_{n_x} \end{bmatrix} 
    H_u
    \begin{bmatrix} Q_1 & \\ & \bm{I}_{n_x} \end{bmatrix} 
    = \begin{bmatrix} V_1 & V_2 & \\ &  & \\ \tilde{M}_2^1 & \tilde{M}_2^2 & M_3 \end{bmatrix}
\end{equation}
where $\begin{bmatrix} (\tilde{M}_2)_1 &(\tilde{M}_2)_2 \end{bmatrix} = M_2 Q_1$ with $(\tilde{M}_2)_1 \in \mathbb{R}^{n_x \times \rho}$ and $(\tilde{M}_2)_2 \in \mathbb{R}^{n_x \times n_x - \rho}$.

Continuing gaussian elimination, we can obtain
\begin{align}
    \hat{L} &\coloneqq \begin{bmatrix} K_1 & & \\ K_2 & \bm{I}_{n_g-\rho} & \\ K_3 & & \bm{I}_{n_x} \end{bmatrix}
    & &\text{such that} &
    \hat{L}^{-1} &= \begin{bmatrix} K_1^{-1} & & \\ -K_2 K_1^{-1} & \bm{I}_{n_g-\rho} & \\ -K_3 K_1^{-1} & & \bm{I}_{n_x} \end{bmatrix}
\end{align}
using $K_3 \coloneqq (\tilde{M}_2)_1 V_1^{-1}$, which leads to
\begin{equation}
    \hat{L}^{-1} 
    \begin{bmatrix} P_1 & \\ & \bm{I}_{n_x} \end{bmatrix} 
    H_u
    \begin{bmatrix} Q_1 & \\ & \bm{I}_{n_x} \end{bmatrix} 
    = \begin{bmatrix} V_1 & V_2 & \\ &  & \\ & K_4 & M_3 \end{bmatrix}
\end{equation}
where $K_4 \coloneqq (\tilde{M}_2)_2 - K_3 V_2$. Because of the zero-row, we must perform row pivoting in order to continue the decomposition. Hence, define
\begin{equation}
    \tilde{P} \coloneqq \begin{bmatrix} \bm{I}_\rho & & \\ & & \bm{I}_{n_x}\\ & \bm{I}_{n_g-\rho} & \end{bmatrix}
\end{equation}
leading to
\begin{equation}
    \tilde{P}
    \hat{L}^{-1}
    \tilde{P}^\top \tilde{P} 
    \begin{bmatrix} P_1 & \\ & \bm{I}_{n_x} \end{bmatrix} 
    H_u
    \begin{bmatrix} Q_1 & \\ & \bm{I}_{n_x} \end{bmatrix} 
    = \begin{bmatrix} V_1 & V_2 & \\ & K_4 & M_3\\ &  & \\\end{bmatrix}
\end{equation}
using the fact that $\tilde{P}^\top \tilde{P} = \bm{I}_{n_x + n_g}$. Now, a remaining LU-decomposition should be performed on the matrix $\begin{bmatrix} K_4 & M_3\end{bmatrix}$ such that
\begin{equation}
    P_2 \begin{bmatrix} K_4 & M_3\end{bmatrix} Q_2 = L_2 \begin{bmatrix} V_3 & V_4 \end{bmatrix}
\end{equation}
where $V_3 \in \mathbb{R}^{n_x \times n_u - \rho}$ and $V_4 \in \mathbb{R}^{n_x \times n_x}$. Note that if $\rho = n_u$, this decomposition reduces to a decomposition of $M_3$. Note that the permutation matrices $P_2$ and $Q_2$ affect also $K_3$ and $V_2$, hence we define
\begin{align}
    \tilde{K}_3 &\coloneqq P_2 K_3\\
    \tilde{V}_2 &\coloneqq V_2 \begin{bmatrix} (Q_2)_{11} & (Q_2)_{12} \end{bmatrix}
\end{align}

Finally, the decomposition of $M$ is given by 
\begin{equation}
    P H_u Q = 
    \begin{bmatrix} K_1 & & \\ \tilde{K}_3 & L_2 & \\ K_2 &  & \bm{I}_{n_g-\rho} \end{bmatrix}
    \begin{bNiceMatrix}[margin=8pt, cell-space-limits=2pt]
        V_1 & \Block[draw=black]{1-2}{\tilde{V}_2} \\
        & V_3 & V_4 \\
        & & \bm{0}_{n_g-\rho \times n_x}
    \end{bNiceMatrix}
    \label{eq:final_decomposition}
\end{equation}
and 
\begin{align}
    P &= 
    \begin{bmatrix} \bm{I}_\rho & & \\ & P_2 & \\ & & \bm{I}_{n_g-\rho} \end{bmatrix}
    \begin{bmatrix} \bm{I}_\rho & & \\ & & \bm{I}_{n_x} \\ & \bm{I}_{n_g-\rho} & \end{bmatrix}
    \begin{bmatrix} P_1 & \\ & \bm{I}_{n_x} \end{bmatrix}\\
    Q &= \begin{bmatrix} Q_1 & \\ & \bm{I}_{n_x}\end{bmatrix} \begin{bmatrix}\bm{I}_\rho &\\ & Q_2 \end{bmatrix}
\end{align}

\begin{algorithm}
    \caption{Blocked LU decomposition of lower triangular matrix $H_u$}
    \begin{algorithmic}
    \Require Matrix $H_u$ given by \eqref{eq:Hu}
    \Ensure LU decomposition with complete pivoting $P H_u P = L U$
    \State Compute $P_1$, $Q_1$, $K_1$, $K_2$, $V_1$ and $V_2$ such that $P_1 M_1 Q_1 = \begin{bmatrix}K1 & \\ K_2 & \bm{I}\end{bmatrix} \begin{bmatrix} V_1 & V_2\\ & \bm{0}\end{bmatrix}$
    \State $\tilde{M}_2 \gets M_2 Q_1$
    \State $K_3 \gets \tilde{M}_2 \begin{bmatrix} \bm{I}_\rho \\ \bm{0}_{n_g-\rho \times \rho} \end{bmatrix} V_1^{-1}$
    \State $K_4 \gets \tilde{M}_2 \begin{bmatrix} \bm{0}_{\rho \times n_g-\rho} \\ \bm{I}_{n_g-\rho} \end{bmatrix} - K_3 V_2$
    \State Compute $P_2$, $Q_2$, $L_2$, $V_3$ and $V_4$ such that $P_2 \begin{bmatrix} K_4 & M_3 \end{bmatrix} Q_2 = L_2 \begin{bmatrix} V_3 & V_4 \end{bmatrix}$
    \State $\tilde{K}_3 \gets P_2 K_3$
    \State $\tilde{V}_2 \gets V_2 \begin{bmatrix} (P_2)_{11} & (Q_2)_{12} \end{bmatrix}$
    \State $L \gets \begin{bmatrix} K_1 & & \\ \tilde{K}_3 & L_2 & \\ K_2 &  & \bm{I}_{n_g-\rho} \end{bmatrix}$
    \State $U \gets \begin{bNiceMatrix}[margin=8pt, cell-space-limits=2pt] V_1 & \Block[draw=black]{1-2}{\tilde{V}_2}\\ & V_3 & V_4 \\ & & \bm{0}_{n_g-\rho \times n_x} \end{bNiceMatrix}$
    \State $P \gets \begin{bmatrix} \bm{I}_\rho & & \\ & P_2 & \\ & & \bm{I}_{n_g-\rho} \end{bmatrix} \begin{bmatrix} \bm{I}_\rho & & \\ & & \bm{I}_{n_x} \\ & \bm{I}_{n_g-\rho} & \end{bmatrix} \begin{bmatrix} P_1 & \\ & \bm{I}_{n_x} \end{bmatrix}$
    \State $Q \gets \begin{bmatrix} Q_1 & \\ & \bm{I}_{n_x}\end{bmatrix} \begin{bmatrix}\bm{I}_\rho &\\ & Q_2 \end{bmatrix}$

    \end{algorithmic}
\end{algorithm}











\section{Blocked Upper Triangular LU decomposition}
The structure in the matrix given by \eqref{eq:Hu-upper} can be exploited by performing the decomposition more efficiently. More precisely, we wish to compute
\begin{equation}
    P H_u Q = L U
\end{equation}
Suppose $M_1$ is decomposed as 
\begin{align}
    P_1 M_1 Q_1 &= L_1 U_1
\end{align}
where
\begin{align}
    L_1 &= \begin{bmatrix} K_1 & \\ K_2 & \bm{I}_{n_x-\rho} \end{bmatrix} & 
    &\text{and} &
    U_1 &= \begin{bmatrix} V_1 & V_2 \\ & \bm{0}_{n_x-\rho \times n_x-\rho} \end{bmatrix}
\end{align}
Define 
\begin{align}
    \tilde{L} &\coloneqq \begin{bmatrix} K_1 & & \\ K_2 & \bm{I}_{n_x-\rho} & \\ & & \bm{I}_{n_g} \end{bmatrix}
    & &\text{such that} &
    \tilde{L}^{-1} &= \begin{bmatrix} K_1^{-1} & & \\ -K_2 K_1^{-1} & \bm{I}_{n_x-\rho} & \\ & & \bm{I}_{n_u} \end{bmatrix}
\end{align}
which leads to
\begin{equation}
    \tilde{L}^{-1} 
    \begin{bmatrix} P_1 & \\ & \bm{I}_{n_g} \end{bmatrix} 
    H_u
    \begin{bmatrix} Q_1 & \\ & \bm{I}_{n_u} \end{bmatrix} 
    = \begin{bmatrix} V_1 & V_2 & K_1^{-1}(\tilde{M}_2)_1\\ &  & (\tilde{M}_2)_2 - K_2 K_1^{-1} (\tilde{M}_2)_1\\ & & M_3 \end{bmatrix}
\end{equation}
where $\begin{bmatrix} (\tilde{M}_2)_1 &(\tilde{M}_2)_2 \end{bmatrix} = P_1 M_2$ with $(\tilde{M}_2)_1 \in \mathbb{R}^{\rho \times n_u}$ and $(\tilde{M}_2)_2 \in \mathbb{R}^{n_x - \rho \times n_u}$.

Define $K_3 \coloneqq K_1^{-1} (\tilde{M}_2)_1$ and $K_4 \coloneqq (\tilde{M}_2)_2 - K_2 K_3$ and 
\begin{equation}
    \tilde{Q} \coloneqq \begin{bmatrix} \bm{I}_\rho & & \\ & & \bm{I}_{n_u}\\ & \bm{I}_{n_x-\rho} & \end{bmatrix}
\end{equation}
leading to
\begin{equation}
    \hat{L}^{-1}
    \begin{bmatrix} P_1 & \\ & \bm{I}_{n_x} \end{bmatrix} 
    H_u
    \begin{bmatrix} Q_1 & \\ & \bm{I}_{n_x} \end{bmatrix} 
    \tilde{Q}
    = \begin{bmatrix} V_1 & K_3 & V_2 \\ & K_4 & \\ & M_3 & \\\end{bmatrix}
\end{equation}
Now, a remaining LU-decomposition should be performed on the matrix $\begin{bmatrix} K_4 \\ M_3\end{bmatrix}$ such that
\begin{equation}
    P_2 \begin{bmatrix} K_4 \\ M_3\end{bmatrix} Q_2 = L_2 U_2
\end{equation}
Note that the permutation matrices $P_2$ and $Q_2$ affect also $K_2$ and $K_3$, hence we define
\begin{align}
    \tilde{K}_2 &\coloneqq P_2 \begin{bmatrix} K_2\\ \bm{0}_{n_g \times \rho} \end{bmatrix}\\
    V_3 &\coloneqq K_3 Q_2
\end{align}

Finally, the decomposition of $M$ is given by 
\begin{equation}
    P H_u Q = 
    \begin{bmatrix} K_1 & \\ \tilde{K}_2 & L_2 \end{bmatrix}
    \begin{bmatrix}
        V_1 & V_3 & V_2 \\
        & U_2 & 
    \end{bmatrix}
    \label{eq:final_decomposition-upper}
\end{equation}
and 
\begin{align}
    P &= 
    \begin{bmatrix} \bm{I}_\rho & \\ & P_2 \end{bmatrix}
    \begin{bmatrix} P_1 & \\ & \bm{I}_{n_x} \end{bmatrix}\\
    Q &= 
    \begin{bmatrix} Q_1 & \\ & \bm{I}_{n_x}\end{bmatrix} 
    \begin{bmatrix} \bm{I}_\rho & & \\ & & \bm{I}_{n_u} \\ & \bm{I}_{n_x-\rho} & \end{bmatrix}
    \begin{bmatrix}\bm{I}_\rho & & \\ & Q_2 & \\ & & \bm{I}_{n_x-\rho}\end{bmatrix}
\end{align}

\begin{algorithm}
    \caption{Blocked LU decomposition of upper triangular matrix $H_u$}
    \begin{algorithmic}
    \Require Matrix $H_u$ given by \eqref{eq:Hu-upper}
    \Ensure LU decomposition with complete pivoting $P H_u P = L U$
    \State Compute $P_1$, $Q_1$, $K_1$, $K_2$, $V_1$ and $V_2$ such that $P_1 M_1 Q_1 = \begin{bmatrix}K1 & \\ K_2 & \bm{I}\end{bmatrix} \begin{bmatrix} V_1 & V_2\\ & \bm{0}\end{bmatrix}$
    \State $\tilde{M}_2 \gets P_1 M_2$
    \State $K_3 \gets K_1^{-1} (\tilde{M}_2)_1$
    \State $K_4 \gets (\tilde{M}_2)_2 - K_2 K_3$
    \State Compute $P_2$, $Q_2$, $L_2$, $V_3$ and $V_4$ such that $P_2 \begin{bmatrix} K_4 \\  M_3 \end{bmatrix} Q_2 = L_2 U_2$
    \State $\tilde{K}_2 \gets P_2 \begin{bmatrix} K_3 \\ \bm{0} \end{bmatrix}$
    \State $V_3 \gets K_3 Q_2$
    \State $L \gets \begin{bmatrix} K_1 & \\ \tilde{K}_2 & L_2\end{bmatrix}$
    \State $U \gets \begin{bmatrix} V_1 & V_3 & V_2 \\ & U_2 & \end{bmatrix}$
    \State $P \gets \begin{bmatrix} \bm{I}_\rho & \\ & P_2 \end{bmatrix} \begin{bmatrix} P_1 & \\ & \bm{I}_{n_g} \end{bmatrix}$
    \State $Q \gets \begin{bmatrix} Q_1 & \\ & \bm{I}_{n_x}\end{bmatrix} \begin{bmatrix} \bm{I}_\rho & & \\ & & \bm{I}_{n_u} \\ & \bm{I}_{n_x-\rho} & \end{bmatrix} \begin{bmatrix}\bm{I}_\rho & & \\ & Q_2 & \\ & & \bm{I}_{n_x-\rho} \end{bmatrix}$

    \end{algorithmic}
\end{algorithm}











\subsection{Verification of blocked LU decomposition}
To verify \eqref{eq:final_decomposition-upper}, compute both the left and right-hand side. The right-hand side is given by
\begin{equation}
    \begin{bmatrix} K_1 & \\ \tilde{K}_2 & L_2 \end{bmatrix}
    \begin{bmatrix} V_1 & V_3 & V_2 \\ & U_2 & \end{bmatrix}
    = 
    \begin{bmatrix} K_1 V_1 & K_1 V_3 & K_1 V_2 \\ \tilde{K}_2 V_1 & \tilde{K}_2 V_3 + L_2 U_2 & \tilde{K}_2 V_2 \end{bmatrix}
    \label{eq:right-hand-side-upper}
\end{equation}
where
\begin{align}
    L_2 U_2 &= P_2 \begin{bmatrix} K_4 \\ M_3 \end{bmatrix} Q_2\\
    &= P_2 \begin{bmatrix} (\tilde{M}_2)_2 - K_2 K_3 \\ M_3 \end{bmatrix} Q_2\\
    &= P_2 \begin{bmatrix} (\tilde{M}_2)_2 \\ M_3 \end{bmatrix} Q_2 - P_2 \begin{bmatrix} K_2 \\ \bm{0} \end{bmatrix} K_3 Q_2\\
    &= P_2 \begin{bmatrix} (\tilde{M}_2)_2 \\ M_3 \end{bmatrix} Q_2 - \tilde{K}_2 V_3
\end{align}

The left-hand side can be using the definitions
\begin{align}
    X_1 \coloneqq \begin{bmatrix} P_1 & \\ & \bm{I}_{n_x} \end{bmatrix}
    H_u
    \begin{bmatrix} Q_1 & \\ & \bm{I}_{n_x}\end{bmatrix}
    &= \begin{bmatrix} P_1 M_1 Q_1 & P_1 M_2 \\ & M_3 \end{bmatrix}\\
    &= \begin{bmatrix} L_1 U_1 & \tilde{M}_2 \\ & M_3 \end{bmatrix}\\
    &= \begin{bmatrix} K_1 V_1 & K_1 V_2 & (\tilde{M}_2)_1 \\ K_2 V_1 & K_2 V_2 & (\tilde{M}_2)_2 \\ & & M_3 \end{bmatrix}\\
    X_2 \coloneqq X_1 \begin{bmatrix} \bm{I}_\rho & & \\ & & \bm{I}_{n_u} \\ & \bm{I}_{n_x-\rho} & \end{bmatrix}
    &= \begin{bmatrix} K_1 V_1 & (\tilde{M}_2)_1 & K_1 V_2 \\ K_2 V_1 & (\tilde{M}_2)_2 & K_2 V_2 \\ & M_3 & \end{bmatrix}
\end{align}
Such that we can compute
\begin{align}
    P H_u Q &= \begin{bmatrix} \bm{I}_\rho & \\ & P_2\end{bmatrix}
    X_2
    \begin{bmatrix}\bm{I}_\rho & & \\ & Q_2 & \\ & & \bm{I}_{n_x-\rho} \end{bmatrix}\\
    &= \begin{bmatrix} K_1 V_1 & (\tilde{M}_2)_1 & K_1 V_2 \\ \tilde{K}_2 V_1 & P_2 \begin{bmatrix} (\tilde{M}_2)_2\\ M_3 \end{bmatrix} & \tilde{K}_2 V_2 \end{bmatrix}
    \begin{bmatrix}\bm{I}_\rho & & \\ & Q_2 & \\ & & \bm{I}_{n_x-\rho} \end{bmatrix}\\
    &= \begin{bmatrix} K_1 V_1 & K_1 V_3 & K_1 V_2 \\ \tilde{K}_2 V_1 & P_2 \begin{bmatrix} (\tilde{M}_2)_2\\ M_3 \end{bmatrix} Q_2 & \tilde{K}_2 V_2 \end{bmatrix}
    \label{eq:left-hand-side-upper}
\end{align}











\section{carry-over of constraints}
In some cases, constraints on $x_{k+1}$ are transferred to stage $k$ using the dynamics equations. In this case, that leads to an augmented matrix 
\begin{equation}
    H_u = \begin{bNiceMatrix}[first-row, first-col]
        & n_u & n_x\\
        n_g & M_1 &    \\ 
        n_x & M_2 & M_3\\ 
        n_c & & M_4
    \end{bNiceMatrix} \label{eq:Hu-carry}
\end{equation}
Note that this is still a lower triangular matrix. If the definition of \eqref{eq:Hu-upper} was used, the structure would now be broken. Additionally, because the dynamics are given by $x_{k+1} = z_k$, the constraints are transferred completely to the columns corresponding to $x_k$. The same approach as explained above can be used to compute the decomposition of this matrix by considering $M_2' \coloneqq \begin{bmatrix} M_2^\top & \end{bmatrix}^\top$ instead of $M_2$ and $M_3' \coloneqq \begin{bmatrix} M_3^\top & M_4^\top \end{bmatrix}^\top$ instead of $M_3$.


\section{Blocked LU Decomposition with carry over of constraints}
To decompose the matrix in \eqref{eq:Hu-carry}, we first write
\begin{align}
    P_1 M_1 Q_1 &= L_1 U_1
\end{align}
where
\begin{align}
    L_1 &= \begin{bmatrix} K_1 & \\ K_2 & \bm{I}_{n_g-\rho} \end{bmatrix} & 
    &\text{and} &
    U_1 &= \begin{bmatrix} V_1 & V_2 \\ & \bm{0}_{n_g-\rho \times n_u-\rho} \end{bmatrix}
\end{align}
Define 
\begin{align}
    \tilde{L} &\coloneqq \begin{bmatrix} K_1 & & \\ K_2 & \bm{I}_{n_g-\rho} & \\ & & \bm{I}_{n_x+n_c} \end{bmatrix}
    & &\text{such that} &
    \tilde{L}^{-1} &= \begin{bmatrix} K_1^{-1} & & \\ -K_2 K_1^{-1} & \bm{I}_{n_g-\rho} & \\ & & \bm{I}_{n_x+n_c} \end{bmatrix}
\end{align}
which leads to
\begin{equation}
    \tilde{L}^{-1} 
    \begin{bmatrix} P_1 & \\ & \bm{I}_{n_x+n_c} \end{bmatrix} 
    H_u
    \begin{bmatrix} Q_1 & \\ & \bm{I}_{n_x} \end{bmatrix} 
    = \begin{bNiceMatrix}[first-row, first-col, last-col=4]
        & \rho & n_u - \rho & n_x & \\
        & V_1 & V_2 & & \rho \\
        & &  & & n_g-\rho\\ 
        & (\tilde{M}_2)_1 & (\tilde{M}_2)_2 & M_3 & n_x\\ 
        & & & M_4 & n_c
    \end{bNiceMatrix}
\end{equation}
where $\begin{bmatrix} (\tilde{M}_2)_1 &(\tilde{M}_2)_2 \end{bmatrix} = M_2 Q_1$ with $(\tilde{M}_2)_1 \in \mathbb{R}^{n_x \times \rho}$ and $(\tilde{M}_2)_2 \in \mathbb{R}^{n_x \times n_x - \rho}$.

Continuing gaussian elimination, we can obtain
\begin{align}
    \hat{L} &\coloneqq \begin{bmatrix} K_1 & & & \\ K_2 & \bm{I}_{n_g-\rho} & & \\ K_3 & & \bm{I}_{n_x} & \\ & & & \bm{I}_{n_c}\end{bmatrix}
    & &\text{such that} &
    \hat{L}^{-1} &= \begin{bmatrix} K_1^{-1} & & &\\ -K_2 K_1^{-1} & \bm{I}_{n_g-\rho} & &\\ -K_3 K_1^{-1} & & \bm{I}_{n_x}& \\ & & & \bm{I}_{n_c} \end{bmatrix}
\end{align}
using $K_3 \coloneqq (\tilde{M}_2)_1 V_1^{-1}$, which leads to
\begin{equation}
    \hat{L}^{-1} 
    \begin{bmatrix} P_1 & \\ & \bm{I}_{n_x+n_c} \end{bmatrix} 
    H_u
    \begin{bmatrix} Q_1 & \\ & \bm{I}_{n_x} \end{bmatrix} 
    = \begin{bNiceMatrix}[first-row, first-col, last-col=4]
        & \rho & n_u - \rho & n_x & \\
        & V_1 & V_2 & & \rho\\
        & &  & & n_g - \rho\\
        & & K_4 & M_3 & n_x\\
        & & & M_4 & n_c
    \end{bNiceMatrix}
\end{equation}
where $K_4 \coloneqq (\tilde{M}_2)_2 - K_3 V_2$. Because of the zero-row, we must perform row pivoting in order to continue the decomposition. Hence, define
\begin{equation}
    \tilde{P} \coloneqq \begin{bmatrix} \bm{I}_\rho & & \\ & & \bm{I}_{n_x+n_c}\\ & \bm{I}_{n_g-\rho} & \end{bmatrix}
\end{equation}
leading to
\begin{equation}
    \tilde{P}
    \hat{L}^{-1}
    \tilde{P}^\top \tilde{P} 
    \begin{bmatrix} P_1 & \\ & \bm{I}_{n_x+n_c} \end{bmatrix} 
    H_u
    \begin{bmatrix} Q_1 & \\ & \bm{I}_{n_x} \end{bmatrix} 
    = \begin{bNiceMatrix}[first-row, first-col, last-col=4]
        & \rho & n_u - \rho & n_x & \\
        & V_1 & V_2 & & \rho\\
        & & K_4 & M_3 & n_x\\
        & & & M_4 & n_c\\
        & &  & & n_g-\rho\\
    \end{bNiceMatrix}
\end{equation}
using the fact that $\tilde{P}^\top \tilde{P} = \bm{I}_{n_x + n_c + n_g}$. Now, a remaining LU-decomposition should be performed on the matrix $\begin{bmatrix} K_4 & M_3\\ & M_4\end{bmatrix}$ such that
\begin{equation}
    P_2 \begin{bmatrix} K_4 & M_3\\ & M_4\end{bmatrix} Q_2 = L_2 U_2
\end{equation}
Note that this matrix also has structure that can be exploited!

Note also that the permutation matrices $P_2$ and $Q_2$ affect also $K_3$ and $V_2$, hence we define
\begin{align}
    \tilde{K}_3 &\coloneqq P_2 \begin{bmatrix} K_3\\ \bm{0} \end{bmatrix}\\
    \tilde{V}_2 &\coloneqq \begin{bmatrix} V_2 & \bm{0} \end{bmatrix} Q_2
\end{align}

Finally, the decomposition of is given by 
\begin{equation}
    P H_u Q = 
    \begin{bNiceMatrix}[first-row, first-col]
        & \rho & n_x + n_c & n_g-\rho\\ 
        \rho & K_1 & & \\
        n_x + n_c & \tilde{K}_3 & L_2 & \\ 
        n_g-\rho & K_2 &  & \bm{I}_{n_g-\rho}
    \end{bNiceMatrix}
    \begin{bNiceMatrix}[first-row, first-col, last-col=3]
        & \rho & n_u - \rho + n_x & \\
        & V_1 & \tilde{V}_2 & \rho \\
        & & U_2 & n_x\\
        & & \bm{0} & n_u - \rho
    \end{bNiceMatrix}
    \label{eq:final_decomposition-carry}
\end{equation}
and 
\begin{align}
    P &= 
    \begin{bmatrix} \bm{I}_\rho & & \\ & P_2 & \\ & & \bm{I}_{n_g-\rho} \end{bmatrix}
    \begin{bmatrix} \bm{I}_\rho & & \\ & & \bm{I}_{n_x+n_c} \\ & \bm{I}_{n_g-\rho} & \end{bmatrix}
    \begin{bmatrix} P_1 & \\ & \bm{I}_{n_x+n_c} \end{bmatrix}\\
    Q &= \begin{bmatrix} Q_1 & \\ & \bm{I}_{n_x}\end{bmatrix} \begin{bmatrix}\bm{I}_\rho &\\ & Q_2 \end{bmatrix}
\end{align}

\begin{algorithm}
    \caption{Blocked LU decomposition of lower triangular matrix $H_u$ with carry-over}
    \begin{algorithmic}
    \Require Matrix $H_u$ given by \eqref{eq:Hu-carry}
    \Ensure LU decomposition with complete pivoting $P H_u P = L U$
    \State Compute $P_1$, $Q_1$, $K_1$, $K_2$, $V_1$ and $V_2$ such that $P_1 M_1 Q_1 = \begin{bmatrix}K1 & \\ K_2 & \bm{I}\end{bmatrix} \begin{bmatrix} V_1 & V_2\\ & \bm{0}\end{bmatrix}$
    \State $\tilde{M}_2 \gets M_2 Q_1$
    \State $K_3 \gets (\tilde{M}_2)_1 V_1^{-1}$
    \State $K_4 \gets (\tilde{M}_2)_2 - K_3 V_2$
    \State Compute $P_2$, $Q_2$, $L_2$ and $U_2$ such that $P_2 \begin{bmatrix} K_4 & M_3\\ & M_4 \end{bmatrix} Q_2 = L_2 U_2$
    \State $\tilde{K}_3 \gets P_2 \begin{bmatrix} K_3\\ 0 \end{bmatrix}$
    \State $\tilde{V}_2 \gets \begin{bmatrix} V_2 & 0 \end{bmatrix} Q_2$
    \State $L \gets \begin{bmatrix} K_1 & & \\ \tilde{K}_3 & L_2 & \\ K_2 &  & \bm{I}_{n_g-\rho} \end{bmatrix}$
    \State $U \gets \begin{bmatrix} V_1 & \tilde{V}_2\\ & U_2\\ & \end{bmatrix}$
    \State $P \gets \begin{bmatrix} \bm{I}_\rho & & \\ & P_2 & \\ & & \bm{I}_{n_g-\rho} \end{bmatrix} \begin{bmatrix} \bm{I}_\rho & & \\ & & \bm{I}_{n_x+n_c} \\ & \bm{I}_{n_g-\rho} & \end{bmatrix} \begin{bmatrix} P_1 & \\ & \bm{I}_{n_x+n_c} \end{bmatrix}$
    \State $Q \gets \begin{bmatrix} Q_1 & \\ & \bm{I}_{n_x}\end{bmatrix} \begin{bmatrix}\bm{I}_\rho &\\ & Q_2 \end{bmatrix}$

    \end{algorithmic}
\end{algorithm}

\section{Numerical results} \FloatBarrier
Figure \ref{fig:scaling} shows the computation time of the full LU decomposition (direct) and the blocked LU decomposition for varying $n_x$, relative rank of $M_1$, $n_u$, and $n_g$. Figure \ref{fig:scaling_relative} shows the relative computation time between the two methods. For small values of $n_x$, the blocked LU is slower than the full LU decomposition, but for other cases, a speedup of 20\% to 30\% is observed.

Figure \ref{fig:scaling_sliced} shows the relative difference in computation time as a function of $n_x$ and $n_u$ for different values of $n_g$. Again, for low values of $n_x$, the blocked LU is slower than the full LU decomposition. Additionally, the speedup is lower if $n_g$ is low.

Finally, Figure \ref{fig:scaling_sliced_rank} shows the relative difference in computation time as a function of $n_x$ and $n_u$ for different values of the relative rank of $M_1$. Again, for low values of $n_x$, the blocked LU is slower than the full LU decomposition. Additionally, the speedup is lower if the relative rank of $M_1$ is low. The rank itself does not seem to influence the relative computation time significantly.

Both of these observations are to be expected since for $n_g = 0$, $H_u$ simply becomes
\begin{equation}
    H_u = \begin{bmatrix} M_1 & 0\end{bmatrix}
\end{equation} 
If $n_u$ is then also small, both approaches just perform the same decomposition of $M_1$. Similarly, if $n_x$ is low, both approaches just perform the decomposition of $M_3$.


\begin{figure}
    \centering
    \begin{subfigure}{0.24\linewidth}
        \includegraphics[width=\linewidth]{../figures/scaling_nx.png}
    \end{subfigure}
    \begin{subfigure}{0.24\linewidth}
        \includegraphics[width=\linewidth]{../figures/scaling_rank.png}
    \end{subfigure}
    \begin{subfigure}{0.24\linewidth}
        \includegraphics[width=\linewidth]{../figures/scaling_nu.png}
    \end{subfigure}
    \begin{subfigure}{0.24\linewidth}
        \includegraphics[width=\linewidth]{../figures/scaling_ng.png}
    \end{subfigure}
    \caption{Comparison in computation time between full LU decomposition (direct) and blocked LU decomposition.}
    \label{fig:scaling}
\end{figure}

\begin{figure}
    \centering
    \begin{subfigure}{0.24\linewidth}
        \includegraphics[width=\linewidth]{../figures/scaling_relative_nx.png}
    \end{subfigure}
    \begin{subfigure}{0.24\linewidth}
        \includegraphics[width=\linewidth]{../figures/scaling_relative_rank.png}
    \end{subfigure}
    \begin{subfigure}{0.24\linewidth}
        \includegraphics[width=\linewidth]{../figures/scaling_relative_nu.png}
    \end{subfigure}
    \begin{subfigure}{0.24\linewidth}
        \includegraphics[width=\linewidth]{../figures/scaling_relative_ng.png}
    \end{subfigure}
    \caption{Relative computation time between full LU decomposition (direct) and blocked LU decomposition.}
    \label{fig:scaling_relative}
\end{figure}

\begin{figure}
    \centering
    \includegraphics[width=\linewidth]{../figures/scaling_sliced.png}
    \caption{Comparison in computation time between full LU decomposition (direct) and blocked LU decomposition as a function of $n_x$ and $n_u$ for varying ranges of $n_g$. Note that we only consider cases where $n_g \leq n_x + n_u$.}
    \label{fig:scaling_sliced}
\end{figure}

\begin{figure}
    \centering
    \includegraphics[width=\linewidth]{../figures/scaling_sliced_rank.png}
    \caption{Comparison in computation time between full LU decomposition (direct) and blocked LU decomposition as a function of $n_x$ and $n_u$ for varying ranges of the relative rank of $M_1$.}
    \label{fig:scaling_sliced_rank}
\end{figure}





\section{Conclusion}
For some values of $n_x$, $n_u$ and $n_g$, the LU decomposition in the Ricatti recursion takes a significant percentage of the computation time. Hence, speeding it up can reduce the overall computation time. In the case of implicit integrators, the constraint jacobian (of which the LU decomposition is computed) is a 2x2 block matrix of which the top-left matrix is zero which can be exploited.

Numerical results show that the time required to compute the decomposition can be reduced by 20-40\% by using a blocked LU deomposition instead of a direct LU decomposition directly. Even though the blocked LU decomposition is not always faster to compute, the parameter values for which the blocked LU is slower correspond to parameter values where the LU decomposition is less important within the recursion.

\end{document}